We collected \rustc{}'s soundness issues reported between January 1, 2022, and
September 1, 2025, focusing on issues pertinent to recent compiler releases.
%
This timeframe aligns with the two most recent \textit{editions} (boundaries for
introducing backward-incompatible changes) of the language:
%
Rust 2021, released on October 21, 2021, and Rust 2024, released on February 20,
2025.
%
Consequently, our target timeframe encompasses over three and a half years,
corresponding to the active use of the Rust 2021 and 2024 editions.

% Each edition is introduced through a specific compiler release that stabilizes
% new language features.
%
% with \rustc{} 1.56.0, with \rustc{} 1.85.0

Since the issue tracker of the official GitHub repository for \rustc{} contains
various kinds of issues, we filtered for soundness issues through automatic
filtering based on the labels assigned to each issue and manual inspection.
%
We initially collected issues related to type checking by utilizing the area
labels (prefixed with ``A-''), which indicate the issue-relevant compiler area,
component, or language feature.
%
Specifically, we retained issues with the fourteen labels classified as
type-related ones by previous work~\cite{liu2025empiricala}:
%
A-type-system, A-inference, A-closures, A-coercions, A-const-generics, A-DSTs,
A-zst, A-trait-system, A-impl-trait, A-trait-objects, A-auto-traits,
A-implied-bounds, A-coinduction, and A-coherence.
%
This yielded 969 issues.

% To collect soundness bug issues within this timeframe, we scraped issue
% metadata from the Rustc GitHub repository.
%
% For example, the A-trait-system label identifies issues related to the Rust
% trait system.
%
% Using the GitHub REST API, we automatically collected metadata for all issues
% within our target timeframe marked with at least one of these 14 labels.

Since this set still includes issues other than soundness (e.g., bugs in
documentation or non-bugs like feature proposals), we further refined the
dataset using two criteria.
%
First, we kept only issues marked with either C-bug or I-unsound, which
signifies that the issue is a bug report or related to type soundness.
%
Second, we filtered out issues marked with specific labels indicating
irrelevance to soundness, which we categorize as follows:
%
\begin{itemize}[leftmargin=*]
  %
  \item Documentation and tool issues: A-docs, A-rustdoc, A-rustdoc-json,
    A-rustdoc-ui, T-rustdoc, and A-clippy, which represent wrong documentation
    or issues in external tools.
  %
  \item Diagnostics issues: A-diagnostics, A-suggestion-diagnostics,
    NLL-diagnostics, and const-generics-bad-diagnostics, which are pertinent to
    the quality of error messages.
  %
  \item Alternative symptoms: I-crash, I-ICE, I-hang, I-compilemem,
    I-compiletime, and A-codegen, which represent symptoms like compiler crashes
    or performance issues.
  %
  \item Non-bugs: C-tracking-issue, C-feature-request, and C-enhancement, which
    track discussions on certain features or represent feature proposals.
  %
  \item Feature-gated features: labels prefixed with ``F-,'' which designate
    issues relevant only to language features guarded by feature gates.
  %
\end{itemize}
%
This filtering process reduced the initial dataset to 320 issues.

Then, we manually investigated each issue to determine whether it is a soundness
issue based on two criteria.
%
First, we retained an issue only if it contained a Rust code snippet that the
reporter claimed the compiler should reject.
%
Second, we excluded issues where the buggy behavior was not caused by the
compiler, e.g., those caused by breaking changes in the standard library.
%
Through this manual inspection, we identified 27 soundness issues.

Finally, we checked whether each issue was a duplicate of another based on
discussions among \rustc{} developers in the issue tracker.
%
We identified 4 issues as duplicates of others already within our dataset and
subsequently excluded them.
%
Additionally, 2 issues were duplicates of issues reported before 2022;
%
we chose to retain these, as their inclusion does not introduce redundancy
within our specific timeframe (see \cref{tab:issues} for these issues and their
original issue numbers).
%
This resulted in a dataset of 23 unique soundness issues.

% We manually inspected the issue description and examined the authors' claims
% regarding the expected behavior.

After the collection, we compared our dataset with that of Liu et
al.~\cite{liu2025empiricala}, which classifies 22 issues as soundness issues.
%
Our inspection shows that only 4 issues overlap across the two datasets, while
19 are unique to ours and 18 are unique to the previous study.
%
Among the 19 issues unique to ours, 4 fall outside the timeframe of the previous
study, 2 were classified under other bug categories, 3 were treated as
duplicates, and 10 were excluded, although we could not find the reasons for
their exclusion in their paper or artifact.
%
This shows that our study provides a more comprehensive view of soundness
issues, beyond an incremental contribution.
%
Among the 18 issues unique to the previous study, 8 have type-unrelated area
labels, 1 lacks both C-bug and I-unsound labels, 8 have ``F-'' labels, and 1 is
a duplicate.
%
We investigated the 9 issues (8 with type-unrelated labels and 1 with no
C-bug/I-unsound label) and found that 7 satisfy our manual inspection criteria
(a Rust code snippet that is claimed the compiler should reject, and buggy
behavior caused by the compiler).
%
Therefore, we complemented our dataset with these 7 issues, resulting in the
final dataset of 30 issues.
%
See \cref{tab:issues} for the issues that overlap with or were added from the
previous study.
